\begin{table}[!htbp]
\centering
\caption{Result-block provenance and parameter scope.}
\label{tab_result_block_provenance}
\footnotesize
\setlength{\tabcolsep}{3.4pt}
\renewcommand{\arraystretch}{1.12}
\begin{tabularx}{\textwidth}{@{}>{\raggedright\arraybackslash}p{0.25\textwidth}>{\centering\arraybackslash}p{0.11\textwidth}>{\raggedright\arraybackslash}X>{\raggedright\arraybackslash}p{0.24\textwidth}@{}}
\toprule
Block & Density & Parameter scope and source & Use in manuscript \\
\midrule
Table~\ref{tab_core_mechanism} / Fig.~\ref{fig_core_evidence} &
0.15 &
Core diagnostic episodes using 18 km/h target speed and TTC threshold 12.0 s. &
Main artifact and mechanism result. \\
\addlinespace[2pt]
Fig.~\ref{fig_ttc_sensitivity} &
0.15 &
Independent sensitivity rerun with stored speed settings and TTC thresholds 6/8/10/12. &
Trend evidence only. \\
\addlinespace[2pt]
Fig.~\ref{fig_target_speed_sensitivity} &
0.15 &
Independent sensitivity rerun with target speeds 15/18/22 and TTC threshold 12.0 s. &
Secondary trend evidence. \\
\addlinespace[2pt]
Table~\ref{tab_density_summary} / Fig.~\ref{fig_density} &
0.08--0.25 &
Density-stress diagnostics using the stored variant configurations. &
Operating-envelope evidence. \\
\bottomrule
\end{tabularx}
\begin{tablenotes}
\item Rows with the same nominal TTC threshold can differ because sensitivity plots are independent diagnostic reruns, not duplicate evaluations of the core Table~\ref{tab_core_mechanism} rows.
\end{tablenotes}
\end{table}
